\begin{textsample}{POS Dim 3 – df – Score 14.00 – df/df\_inf\_18.txt}  \label{ex:f3_pos_130}
13.
TRANSIÇÕES NA EDUCAÇÃO \textbf{INFANTIL} As transições estão presentes na Educação \textbf{Infantil} das mais diversas formas : transição de casa para a \textbf{instituição} de Educação \textbf{Infantil} ; transição de uma \textbf{instituição} de Educação \textbf{Infantil} para outra, tais como da \textbf{instituição} \textbf{parceira} para a pública ; transição no interior da própria \textbf{instituição} educativa e transição da Educação \textbf{Infantil} para o Ensino Fundamental.
É importante mencionar que a transição de casa para a Educação \textbf{Infantil} pode ocorrer em qualquer período da \textbf{infância}, ou seja, pode ser entre os \textbf{bebês}, as \textbf{crianças} bem \textbf{pequenas} e as \textbf{crianças} \textbf{pequenas}.
Nesse sentido, a atenção ao acolhimento e às estratégias pedagógicas para esse momento precisam considerar as especificidades de cada um nesses períodos, observando as necessidades de cada \textbf{criança}.
É preciso sensibilidade para o acolhimento, para a inserção e para as diversas possibilidades de transição que ocorrem na Educação \textbf{Infantil}, tais como períodos prolongados em que a \textbf{criança} fica afastada da \textbf{instituição} educativa e, ao retornar, \textbf{depara} -se com algum tipo de conflito por estar novamente adentrando um espaço que se diferencia, em vários aspectos, de sua casa ; transições que ocorrem entre os períodos de férias ou de passagem de um ano para outro, entre outras.
A passagem do conhecido para o desconhecido pode desencadear sentimentos de ansiedade, expectativas positivas e negativas, tensões, estresses, medos, traumas e crises, que, caso ocorram, incidem sobre o desenvolvimento integral da \textbf{criança} ( FACCI, 2004 ).
Aos \textbf{adultos} cabe um olhar cuidadoso e uma postura acolhedora e afetuosa sobre os processos vivenciados pela \textbf{criança}, criando estratégias adequadas aos diferentes momentos de acolhida, inserção e transição.
Assim, durante a inserção inicial, as \textbf{instituições} que ofertam Educação \textbf{Infantil} devem favorecer um ambiente físico e social onde as \textbf{crianças} se sintam protegidas, acolhidas e seguras para arriscarem e enfrentarem desafios.
Em relação à transição para o Ensino Fundamental, as DCNEI recomendam : > Na transição para o Ensino Fundamental a proposta pedagógica deve prever formas para garantir a continuidade no processo de aprendizagem e desenvolvimento das \textbf{crianças}, respeitando as especificidades \textbf{etárias}, sem antecipação de conteúdos que serão trabalhados no Ensino Fundamental ( BRASIL, 2010a, p. 30 ).
É perfeitamente possível uma passagem instigante e \textbf{interessante} entre as etapas da Educação Básica.
Ao inserir -se no Ensino Fundamental, não é preciso que os \textbf{pequenos} se \textbf{deparem} com um hiato entre as experiências vivenciadas na Educação \textbf{Infantil} e as práticas educativas da nova etapa.
Nesse sentido, o projeto da Plenarinha tem contribuído ao buscar desenvolver atividades que envolvam tanto as \textbf{crianças} da Educação \textbf{Infantil} como as que se encontram no primeiro ano do Ensino Fundamental.
Evidencia -se, portanto, a necessidade de se estabelecer um diálogo entre as etapas, com ações que superem a tradicional dicotomia que tem contaminado essa passagem.
Seguem algumas \textbf{sugestões} para as \textbf{instituições} de educação coletiva para a primeira \textbf{infância}, visando minimizar os impactos que ocorrem em momentos de transição : - perceber a convergência necessária entre as etapas, tendo a educação como um direito das \textbf{crianças}, compreendendo -as como sujeitos de cultura e cidadãos de direitos ; - ler, estudar e discutir os currículos tanto da Educação \textbf{Infantil} quanto do Ensino Fundamental, mais especificamente dos anos que compreendem o Bloco Inicial de Alfabetização – BIA do 2º Ciclo ; - possibilitar momentos de visita e primeiro contato com a \textbf{instituição} educativa que receberá a \textbf{criança} da Educação \textbf{Infantil} no ano seguinte ; - envolver as famílias e/ou responsáveis no processo de transição entre as etapas, por se tratar de um momento de insegurança e dúvidas para muitos.
Outra questão que merece nota é afirmar que a Educação \textbf{Infantil} não tem por intuito preparar as \textbf{crianças} para o Ensino Fundamental.
É certo que, na condição de etapas da Educação Básica, a Educação \textbf{Infantil} e o Ensino Fundamental precisam estabelecer uma articulação, entendendo que a \textbf{criança} que chega a essa etapa continua sendo \textbf{criança} e precisa ser compreendida dentro de suas especificidades.
De acordo com as DCNEI ( 2010a ) e a BNCC ( 2017 ), a natureza, a identidade e os objetivos de aprendizagem e desenvolvimento presumidos nos dispositivos legais não abordam a alfabetização como uma obrigação na Educação \textbf{Infantil}.
É evidente que, nessa fase, a \textbf{criança} já inicia seu processo de leitura de mundo, por meio de inúmeras atividades, mas isso ocorre de uma forma mais ampla, para além da codificação ou decodificação da língua escrita.
A primeira etapa da Educação Básica tem finalidades próprias que devem ser alcançadas na perspectiva do desenvolvimento \textbf{infantil}, ao se respeitar as \textbf{brincadeiras} e interações e o \textbf{cuidar} e educar, no tempo singular da primeira \textbf{infância}.

% matched lemmas: adulto, bebê, brincadeira, criança, cuidar, deparar, etário, infantil, infância, instituição, interessante, parceiro, pequeno, sugestão
\end{textsample}
